\documentclass[11pt]{article}
\usepackage[margin=0.85in]{geometry}
\usepackage{amsmath,amssymb}
\usepackage{enumitem}
\usepackage{xcolor}
\usepackage{titlesec}
\usepackage{setspace}
\setstretch{1.08}

\definecolor{key}{rgb}{0.10,0.35,0.65}
\titleformat{\section}{\normalfont\large\bfseries\color{key}}{\thesection}{0.6em}{}
\titleformat{\subsection}{\normalfont\normalsize\bfseries}{\thesubsection}{0.5em}{}
\titlespacing{\section}{0pt}{8pt}{4pt}
\titlespacing{\subsection}{0pt}{5pt}{2pt}
\setlist{leftmargin=1.3em,itemsep=1pt,topsep=2pt}
\newcommand{\Q}[1]{\smallskip\noindent\textbf{Q. #1}\par}
\newcommand{\A}[1]{\noindent #1\par}

\pagestyle{empty}

\begin{document}

\begin{center}
  {\Large\bfseries\color{key} Presentation Prep: Adjoint Magnet Shielding} \\[2pt]
  {\normalsize Follow-up check-in, 2026-08-10}
\end{center}

\vspace{-4pt}
\section{What I Am Presenting}
Last meeting the adjoint coil-importance map was validated and the placed-versus-uniform shield test
was proposed with the outcome open. Since then that test resolved, and I turned it into a real
experiment.

The idea behind the whole program is that a single adjoint solve gives a map of how much every point
in the machine contributes to the coil load, for free, without a separate transport run per design.
If that map also told me where to put shield, I could optimize the radial build without running
Monte Carlo in the loop. The experiment asks whether that is true.

There is a ladder of quantities, from cheapest to truest. Attribution says where the neutrons that
load the coil are born. Scalar worth says, to leading order, where blocking flux would help. Signed
worth corrects that for the fact that adding shield means giving up breeder at a fixed envelope.
Measured worth is the real answer: build the geometry, run transport, and see how the coil dose
actually changes. The experiment tests whether any of the cheap rungs predicts the true one.

The measurement covered twelve shield patches. Two results changed the picture. The cheap proxy does
not rank placements, so it cannot be the objective an optimizer follows. It does correctly flag the
one patch where adding shield makes the coil worse, so it works as a guard. And the coil loading is
set by a shared inboard region rather than by twenty independent coil footprints, so magnet
protection is a collective problem and the optimization is low-dimensional.

\section{Numbers To Have Ready}
\begin{itemize}
  \item Peak coil fast flux at equal material: uniform $1.00$, single-coil placement $0.24$
        (a factor of four lower), distributed placement $0.68$.
  \item Rank correlation of measured worth with the proxies (twelve patches): scalar worth $-0.34$,
        attribution $-0.17$, signed worth $+0.16$. None are significant.
  \item Backfire patch: adding shield there raises the coil dose by forty-one percent.
  \item Doses: twelve batches of four million, no weight windows, correlated seed. Coil-20 relative
        error $0.7$ to $1.7$ percent. Dose changes are twenty to a hundred standard deviations.
  \item Attenuation lengths used in the signed worth: tungsten carbide shield about $8$ cm, FLiBe
        breeder about $17$ cm. Breeder floor $15$ cm; breeding ratio floor $1.05$.
  \item Geometry: reactor-scaled quasi-helical build, major radius about $13.7$ m, minor radius
        about $2.2$ m. Reciprocity from last time: quasi-helical $+0.76$, quasi-axisymmetric $+0.83$.
\end{itemize}

\section{The Physics}

\subsection{The Adjoint And The Contributon}
The forward problem is the physical one. Neutrons are born in the plasma with a source $S(r)$, they
transport through roughly a meter of blanket and shield, and a small fraction deposit energy in the
coils. The coil load is $R=\int \Sigma_{\mathrm{resp}}\,\phi\,dV$, the response cross section times
the forward flux, integrated over the coils.

The adjoint problem runs this backward. Put the coil response function in as a source and solve the
transport equation with the operator transposed. The result is the adjoint flux $\psi^\dagger(r)$,
which reads as an importance: the expected contribution to the coil load of one neutron introduced at
$r$. A neutron born deep in the plasma on a clear path to the inboard coil has high importance; one
born where geometry already blocks it has low importance. Because the coil heating response is
direction-blind, the adjoint source is isotropic and a scalar solve is enough for placement.

Reciprocity is the identity $R=\int S\,\psi^\dagger = \int \Sigma_{\mathrm{resp}}\,\phi$. The same
coil load comes out whether computed forward from the source or backward from the importance.
Validating it, per-coil correlation $+0.76$ and $+0.83$ on the two devices, is what says the adjoint
solve is trustworthy.

The contributon is $C(r)=\phi(r)\,\psi^\dagger(r)$. It is the local density of response flow: how
much of the coil load passes through the point $r$. Integrate it over any region and you get that
region's share of the coil load. Everything downstream is the contributon read out in different
places.

\subsection{Attribution Versus Worth}
Attribution is the contributon evaluated at the plasma source, $A=S\,\psi^\dagger$. It answers which
plasma regions cause the coil load. Worth is the contributon summed over the shield band,
$W=\sum_{\mathrm{shield}}\phi\,\psi^\dagger$. It answers where adding shield should intercept the most
response.

They are the same quantity in two places, and the reason to separate them is that they answer
different questions. A plasma region can emit many coil-bound neutrons yet sit on a sightline where
shield cannot easily catch them; a shield location can carry a lot of contributon yet already be
thick enough that more does little. Where the neutrons come from is not the same as where blocking
them helps. The experiment measures how far apart these really are.

\subsection{The Fixed-Envelope Trade}
The outer envelope is fixed. The coils do not move and the machine does not grow, because both cost
far more than shielding does. So shield cannot simply be added. Adding a thickness of shield in a
patch means removing the same thickness of breeder there, and the breeder is the tritium-breeding
blanket. Protecting the magnet costs breeding.

This is why worth must be signed. The benefit is the contributon the added shield intercepts; the
cost is the contributon the removed breeder no longer sees:
$W_s=\sigma_{\mathrm{sh}}\sum_{\mathrm{shield}}C-\sigma_{\mathrm{br}}\sum_{\mathrm{breeder}}C$, with
the $\sigma$ set by the two materials' fast-neutron attenuation. A trade only pays off where the
first term wins.

\subsection{Why The Naive Worth Gets The Sign Wrong}
First-order removal theory says the change in coil load from a small change in cross section is
$\delta R=-\int \delta\Sigma\,\phi\,\psi^\dagger$. For the trade the cross section goes up in the
trade zone, since a centimeter of shield attenuates more than a centimeter of breeder, so this
predicts $\delta R<0$, the dose down, for every patch. It cannot produce a patch where shielding
makes the coil worse.

What it misses is that the coil responds to fast flux, and the breeder does more than remove
neutrons. It moderates them: a fast neutron entering the breeder is likely downscattered on lithium,
beryllium, and fluorine below the energy the coil responds to. Remove breeder and that moderation
goes away, so more neutrons reach the coil still fast. That is a positive contribution to the fast
coil flux that no total-cross-section model contains. On most patches the added shield outweighs it
and the dose falls. On one patch the released fast flux wins and the dose rises by forty-one percent.

Getting the sign right therefore needs a treatment that separates fast from thermal, because the
effect is a transfer between energy groups, not a change in total attenuation. This is why a
single-group signed worth ranks only roughly, and why a fast-group forward flux is the correct
weighting to test next.

\subsection{Why The Problem Is Collective}
The intuitive picture is that each coil has its own line of sight to the plasma, so each coil should
be protected locally. The transport says otherwise. The inboard side of the torus is where the
plasma sits closest to the coils and the flux is highest, and every coil passes near the inboard at
some toroidal angle. The inboard is a hotspot shared by the whole coil set.

Concentrating shield on the inboard band, symmetric across field periods, therefore protects every
coil's inboard segment at once. Spreading the same material by per-coil importance dilutes it and
protects no coil well. Single-coil placement over-protected one coil's inboard so hard that the peak
load jumped to a different, unprotected coil. The lesson is to treat the inboard band as one target
rather than twenty coils, which also makes the optimization low-dimensional.

\section{Anticipated Questions And Answers}
\Q{Is the ranking just Monte Carlo noise?}
\A{No. The dose changes are twenty to a hundred standard deviations for the strong patches, and the
correlated seed cancels most of the variance in the difference. The near-flat worth across most
patches is a real feature, not scatter. The one backfire patch is far outside noise.}

\Q{Why no weight windows on a deep coil dose?}
\A{FW-CADIS weight windows over-split on this geometry, turning two million primaries into tens of
millions of secondary tracks, so the first batch never finishes. Brute-force sampling resolves the
deep coil dose to about one percent and, with a fixed seed, makes each patch dose correlated with the
baseline. The difference is then clean even though the absolute error is modest.}

\Q{Are the signed-worth cross sections tuned to get the answer?}
\A{No. They are the fast-neutron attenuation lengths of the shield and breeder materials. The
four-corner result held across a factor of ten in their ratio, and the twelve-patch result stays weak
across the same sweep, so the conclusion does not depend on the choice.}

\Q{If the proxy does not rank, what is the adjoint good for?}
\A{Three things. It tells you where the coil load comes from, which is physical insight. The signed
worth flags patches where the trade backfires, which a shield-the-hotspot heuristic would miss. And
the importance map is concentrated on the inboard, which collapses the optimization to a few
parameters. The adjoint sets the search space and the constraints, not the objective.}

\Q{So there is no automatic optimizer?}
\A{There is, but not the zero-cost one. The workable system is a low-dimensional Bayesian search
scored by real transport, with the signed worth as a guard that keeps it from degrading the magnet.
The value of the free adjoint is to shrink the search and rule out backfires, not to replace the
transport evaluation.}

\Q{Could the multi-coil result be a budget mismatch or an allocator bug?}
\A{The budgets match to the centimeter. The single-coil field concentrates thirty-five centimeters on
the inboard; the distributed field spreads about ten. The tell is the total coil flux: the
concentrated field cuts it by more than half, the distributed field barely moves it. Concentrating on
the shared inboard helps every coil; spreading by per-coil importance dilutes. A concentrated
distributed field is in transport now to confirm.}

\Q{Trading breeder for shield lowers the breeding ratio. Is the trade feasible?}
\A{That is the binding constraint and the reason for the fixed-envelope framing. The trade keeps a
breeder floor and the optimizer holds the breeding ratio above its floor. The backfire result matters
because a careless placement can hurt both the magnet and the breeding at once.}

\Q{How does this connect to the spin-polarized fuel source?}
\A{The polarized source changes where the coil load peaks through the emission direction. The
attribution map is exactly the tool that links emission direction to coil loading. This shielding
work builds the adjoint framework the polarized source will couple into. The plan is two papers, one
on the directional source with the adjoint, one on the adjoint-informed shielding.}

\Q{Is any of this already published?}
\A{The existing radial-build tool does manual sweeps and names automated optimization as future work.
Polarized-source transport in this code was shown by another group this year, so the source itself is
not our claim. The closed-loop breeder-for-shield magnet optimizer at fixed envelope, the backfire
result, and the collective-hotspot reframing appear to be new.}

\Q{Does angular order matter for placement?}
\A{Not for placement. On the real coil through the shield the first-order correction is under one
percent and the peak location is unchanged. It reaches about seventy percent in a deep scattering slab,
so angular order matters for through-shield magnitude, not for where to place shield.}

\Q{What is next, concretely?}
\A{Two runs finish tomorrow: the concentrated distributed placement to test whether it beats
single-coil concentration, and a fast-group signed worth to test whether the coil-relevant energy
range recovers a usable ranking. After that, the low-dimensional optimizer with the backfire guard,
and the reactor-scale coil sets for a standoff sweep.}

\end{document}
